\bookchapter{第六十九章：速查之冊與資源}

\chapterintro{此章者，可獨用之速查冊也。汝記不清某寫法、或須速決之際，徑翻至此。宜視之為「囊中之卡」：常用碼之速查、常見之問、學習之資，一頁不遺。}

\booksection{一、常用碼速查之卡}

\begin{center}
\begin{tabularx}{\linewidth}{|X|X|}
\hline
\textbf{欲作何事} & \textbf{關鍵之碼} \tabularnewline
\hline
頁面私有之態 & @State private var x = 初值 \tabularnewline
跨頁之邏輯 & \texttt{@Observable class VM \{ var x = 0 \}} \tabularnewline
豎列/橫列之佈 & \texttt{VStack \{ \}}/\texttt{HStack \{ \}} \tabularnewline
彈性之空 & \texttt{Spacer()} \tabularnewline
列表 & List \{ ForEach(陣列) \{ ... \} \} \tabularnewline
頁面之跳轉 & \texttt{NavigationStack + NavigationLink} \tabularnewline
底部彈窗 & .sheet(isPresented: \$show) \{ 內容 \} \tabularnewline
闔當前頁 & \texttt{@Environment(\textbackslash{}.\allowbreak dismiss) var dismiss} \tabularnewline
連網取數據 & \texttt{try await URLSession.\allowbreak shared.\allowbreak data(from: url)} \tabularnewline
JSON 之解析 & \texttt{struct X: Codable}+\texttt{JSONDecoder().\allowbreak decode} \tabularnewline
貯小配置 & UserDefaults.standard.set(值, forKey: "鍵") \tabularnewline
貯結構化數據 & \texttt{@Model}+\texttt{modelContainer}+\texttt{@Query} \tabularnewline
深色模式之配 & 用系統顏色（.primary / Color(.systemGray6)），勿寫死 \tabularnewline
行測試 & CmdU \tabularnewline
\hline
\end{tabularx}
\end{center}

\begin{center}\small\color{gray}iOS 開發速查之卡（2026 年 8 月版）\end{center}

\booksection{二、常見之問 FAQ}

Q1：吾無 Mac，可學 iOS 開發乎？

可學其語，難為全之開發。iPad/iPhone 之 Swift Playgrounds，足學第 61 章之全文；然 Xcode 為 macOS 所獨，終須有一 Mac。預算有限者，可先租雲 Mac 或購二手 Mac mini 以過渡，學成其語再定所投。

Q2：學 iOS 開發，費幾何？

學習之期：0 元（Xcode 免費 + 免費 Apple ID 真機調試）。發布之期：開發者賬號 99 美元/年（約 700 元）。加一 Mac 硬件之投，此其全費也。較其所償——App Store 之上，雖一小具亦有付費之戶——此賬合算。

Q3：SwiftUI 與 UIKit，孰先學？

SwiftUI。2026 年新之項、新之教、新之職，皆以 SwiftUI 為主，本書終始用之。UIKit 為前代之框，猶多用於舊項——待汝 SwiftUI 嫻熟，按需補 UIKit（官方有互通之教），上手甚速。序勿倒置。

Q4：iOS 開發，覓職易乎？

2026 年之市情：徒「能繪界面」之初職，競爭烈；然「曉 Swift + SwiftUI + 有上架之作」之開發者，猶為所爭。iOS 職之俸中數，於諸技之列常居前。學習期之汝，至善之履歷，即商店中自作之 App。

Q5：App 為審核所拒，奈何？

勿惶，此常事也——首提見拒之率不低。蘋果拒之由，常書之甚具（何頁、何條之則），逐條改之，於 App Store Connect 復而說明，再提即可。常見拒之三大由：權限之用不明、界面有 bug、元數據（截圖/描述）與實物不符。先自察此三，可省一輪之候。

\booksection{三、Xcode 高頻快捷鍵速查}

\begin{center}
\begin{tabularx}{\linewidth}{|X|X|}
\hline
\textbf{操作} & \textbf{快捷鍵} \tabularnewline
\hline
運行 / 測試 & CmdR / CmdU \tabularnewline
啟畫布預覽 & OptionCmd \tabularnewline
速啟檔案 & CmdO \tabularnewline
全局搜索 & CmdF \tabularnewline
格式化縮進 & CtrlI \tabularnewline
註釋 / 去註釋 & Cmd/ \tabularnewline
滌構建之緩存 & CmdK \tabularnewline
跳至其定義 & Cmd 點按 \tabularnewline
易名（全局安改） & 右鍵 Refactor Rename \tabularnewline
\hline
\end{tabularx}
\end{center}

\begin{center}\small\color{gray}Xcode 高頻快捷鍵速查\end{center}

\begin{tcbinfo}{小貼士：編譯玄而莫測，先滌緩存}
「昨日可行今日不可」之靈異，什六七為構建緩存之亂。CmdK 滌之而重構，可愈玄惑之大半。不驗則重啟 Xcode（並 Mac 同之），再不驗則「刪 DerivedData 之夾」——此 iOS 開發者口耳相傳之終極三連也。

\end{tcbinfo}

\booksection{四、薦學習之資}

\booksection{六、薦學習之資}

\begin{center}
\begin{tabularx}{\linewidth}{|X|X|X|}
\hline
\textbf{類型} & \textbf{資源} & \textbf{薦之之由} \tabularnewline
\hline
官方 & developer.apple.com & 蘋果官方之檔，SwiftUI/Swift 之教悉在其中，權威居首 \tabularnewline
官方 & Apple 之 SwiftUI Tutorials & 官方親授之程，於新手至為友善 \tabularnewline
官方 & Swift 官網 (swift.org) & 語言官方之檔與在線演練之場 \tabularnewline
社群 & Hacking with Swift & Paul Hudson 免費教程之庫，其質公認居首列 \tabularnewline
社群 & Stack Overflow (swiftui 標籤) & 報錯求救之首驛 \tabularnewline
視頻 & Apple Developer 官方之頻道（WWDC 視頻） & 歲歲新技之權威講解 \tabularnewline
設計 & Apple 人機界面指南（HIG） & 造「蘋果味」界面之官方聖經 \tabularnewline
配套 & 第七篇《AI編程實戰篇》（第37-47章） & 以 AI 輔助 iOS 開發之完整心法 \tabularnewline
\hline
\end{tabularx}
\end{center}

\begin{center}\small\color{gray}iOS 開發學習資源清單（2026 年 8 月）\end{center}

\booksection{五、結語：從第一個 App 到第 N 個 App}

回顧既往，爾自「不知 SwiftUI 為何物」，至裝好 Xcode、學會 Swift、以 SwiftUI 繪出界面、為 App 接上數據庫與網絡，再到測試、打包、上架——此路之上，爾已走完從零到一最難之一段。

技藝之世界變化極快：框架會升級，工具會換代，AI 會接手越來越多「敲鍵盤」之活。然有一物不會過時：爾把模糊之想法，拆解成清晰之步驟，再一步步實現之能力。工具者，放大器也；思路者，地基也。放大器可替換，地基不可動搖。

汝不復徒為"玩手機之人"。汝能以胸中之念，化為數億 iPhone 上實運行之 App。自今日始，以鍵盤創物——此 iOS 開發者之浪漫也。

\booksection{參考文獻}

\begin{enumerate}
  \item Apple Inc. (2026). SwiftUI Framework Documentation. developer.apple.com.
  \item Apple Inc. (2026). Swift Programming Language Guide. swift.org.
  \item Apple Inc. (2026). SwiftData Framework Documentation. developer.apple.com.
  \item Apple Inc. (2026). App Store Review Guidelines. developer.apple.com/app-store/review.
  \item Apple Inc. (2026). Human Interface Guidelines (HIG). developer.apple.com/design.
  \item Apple Inc. (2026). Xcode 27 Release Notes and App Distribution Guide.
\end{enumerate}
